% =====================
\section{スコープ・マネジメント}

\subsection{成果物}

本プロジェクトの最終的な成果物は，
サーボモータで駆動する指揮者人形を各楽器ユニットがセンサーで観測することにより，
楽器間の無線・有線通信なしで自律的に同期し，
「かえるのうた」の5パート輪唱を演奏する\textbf{Arduinoオーケストラシステム}である．

主な構成要素は以下のとおりである．
\begin{itemize}
  \item 親機（指揮者人形ユニット）：Arduino Uno R4 WiFi $\times$ 1台
  \item 子機（楽器ユニット）：Arduino Uno R4 WiFi $\times$ 5台，各PCにシリアル接続
  \item 音響合成：Processing（Minimライブラリ）$\times$ 5台
\end{itemize}

\subsection{プロジェクトの対象範囲と非対象範囲}
% 全員で記載予定
（全員で後から記載）

\subsection{機能分解構造FBS}

機能分解構造（Function Breakdown Structure：FBS）を図\ref{fig:fbs}に示す．
FBSはユーザ視点でシステムが有すべき機能を階層的に分解したものである．

\begin{figure}[H]
  \centering
  \fbox{\begin{minipage}{0.92\linewidth}
    \centering\vspace{4cm}
    \textit{（FBS図を後から挿入：\texttt{base\_plan/fbs.drawio}）}
    \vspace{4cm}
  \end{minipage}}
  \caption{機能分解構造FBS}
  \label{fig:fbs}
\end{figure}

FBSは5つの主機能グループ（L1）とその下位機能（L2）で構成される．
担当B/Cが担当する機能を以下に示す．

\subsubsection*{担当B：指揮動作検出・演奏制御・状態管理}
\begin{description}
  \item[拍タイミング検出（F2a1）]
    フォトトランジスタによって指揮者人形腕先端のLEDを検出し，
    拍の基準時刻を割り込み処理で取得する．
  \item[テンポ推定（F2a3）]
    連続する拍間隔の時系列データから指数移動平均（EMA）によってテンポを推定する．
    通常演奏中は$\alpha=0.2$（緩やかな追従），テンポ変化中は$\alpha=0.5$（素早い追従）を用いる．
  \item[フェルマータ検出（F2a4）]
    赤外線距離センサーによって腕の静止状態を検出し，フェルマータを識別する．
  \item[小節同期・見逃し補完（F2a5）]
    拍が3回連続して未検出の場合に演奏状態を停止へ遷移させるなど，
    センサー取りこぼしへのフォールバック処理を行う．
  \item[演奏状態遷移管理（F2e1）]
    停止・キャリブレーション・予備振り・通常演奏・テンポ変化中・フェルマータ中の
    6状態を管理する有限状態機械を実装する．
  \item[異常時フォールバック（F2e2）]
    3拍連続未検出など異常時に演奏を安全に停止させる処理を担う．
\end{description}

\subsubsection*{担当C：演奏制御（楽譜・通信）}
\begin{description}
  \item[輪唱オフセット制御（F2c1）]
    各楽器ユニットの演奏開始小節を管理し，
    ハイハット・バスから始まり順次旋律パートが参加する輪唱構成を実現する．
  \item[楽譜進行制御（F2c2）]
    小節カウンタに基づいて楽譜配列のポインタを進め，
    テンポ区間（80/104\,bpm）の切り替えを行う．
  \item[テンポ追従演奏（F2c3）]
    担当Bが推定したテンポに追従して，
    NOTE\_ON/NOTE\_OFFパケットを適切なタイミングでProcessingへ送信する．
\end{description}

\subsection{製品分解構造PBS}

製品分解構造（Product Breakdown Structure：PBS）を図\ref{fig:pbs}に示す．
PBSはシステムを構成する成果物を階層的に分解したものである．

\begin{figure}[H]
  \centering
  \fbox{\begin{minipage}{0.92\linewidth}
    \centering\vspace{4cm}
    \textit{（PBS図を後から挿入：\texttt{base\_plan/pbs.drawio}）}
    \vspace{4cm}
  \end{minipage}}
  \caption{製品分解構造PBS}
  \label{fig:pbs}
\end{figure}

担当B/Cが担当する成果物を以下に示す．

\subsubsection*{担当B：子機ユニット（PB）}
\begin{description}
  \item[赤外線センサー配線（PB1a）]
    腕の位置を連続取得するための赤外線距離センサー回路．
    出力電圧$V$から距離$d$\,[cm]を次式で算出する．
    \begin{equation}
      d = 27.728 \times V^{-1.2045}
    \end{equation}
    静止判定・フェルマータ検出・テンポ変化検出に使用する．
  \item[フォトトランジスタ配線（PB1b）]
    指揮者人形腕先端LEDを検出し，デジタル割り込みで拍の基準時刻を取得するための受光回路．
  \item[観測ソフトウェア（PB2a）]
    センサー値の取得・フィルタリング・距離変換処理を実装するArduinoプログラム．
  \item[状態機械ソフトウェア（PB2b）]
    6状態の状態遷移およびEMAによるテンポ推定を実装するArduinoプログラム．
\end{description}

\subsubsection*{担当C：楽譜・通信ユニット（PC）}
\begin{description}
  \item[楽譜配列（PC1a）]
    「かえるのうた」の5パート分の音符データを配列として定義したデータファイル．
  \item[輪唱オフセット定義（PC1b）]
    各楽器の演奏開始小節（0, 0, 2, 4, 6小節目）を定数として定義する．
  \item[遅延補正値（PC1c）]
    Arduino→Processing間のシリアル通信遅延を補正するオフセット値．
    結合テストC1にて実測・確定する．
  \item[演奏制御ソフトウェア（PC2a）]
    楽譜ポインタ管理・NOTE\_ON/OFF送信・フェルマータフラグ制御を実装するArduinoプログラム．
\end{description}

\subsection{FBSとPBSの対応関係}

FBSのすべての機能（L2）に対してPBSのサブ成果物（L3）が1つ以上対応していることを確認する．
担当B/Cに関わる対応関係を表\ref{tab:fbs-pbs}に示す．

\begin{table}[H]
  \centering
  \begin{tabularx}{\linewidth}{llX}
    \toprule
    FBS機能 & PBS成果物 & 対応の説明 \\
    \midrule
    拍タイミング検出（F2a1） &
      フォトTr配線（PB1b）・観測SW（PB2a） &
      フォトTrがLEDを検出し，観測SWが割り込みで拍時刻を記録する \\
    テンポ推定（F2a3） &
      状態機械SW（PB2b） &
      EMAアルゴリズムを状態機械内で実装する \\
    フェルマータ検出（F2a4） &
      赤外線センサー配線（PB1a）・観測SW（PB2a） &
      赤外線センサーで腕の静止状態を検出する \\
    小節同期・見逃し補完（F2a5） &
      状態機械SW（PB2b） &
      3拍連続未検出で停止状態へ遷移するフォールバック処理 \\
    輪唱オフセット制御（F2c1） &
      輪唱オフセット定義（PC1b）・演奏制御SW（PC2a） &
      オフセット定数に基づき演奏開始小節を制御する \\
    楽譜進行制御（F2c2） &
      楽譜配列（PC1a）・演奏制御SW（PC2a） &
      楽譜配列を小節カウンタで進め，テンポ区間を切り替える \\
    テンポ追従演奏（F2c3） &
      演奏制御SW（PC2a）・遅延補正値（PC1c） &
      テンポに合わせてNOTE\_ON/OFFの送信タイミングを制御する \\
    演奏状態遷移管理（F2e1） &
      状態機械SW（PB2b） &
      6状態の遷移条件・遷移処理を実装する \\
    異常時フォールバック（F2e2） &
      状態機械SW（PB2b） &
      異常検出時に演奏を安全に停止させる \\
    \bottomrule
  \end{tabularx}
  \caption{FBSとPBSの対応関係（担当B/C）}
  \label{tab:fbs-pbs}
\end{table}
